\begin{recipe}{Ciabatta}
\label{recipe:ciabatta}
\kicker[Brood,Italiaans,Basisrecept]{Brood \& basis · Italiaans · basisrecept}
\index[register]{Ciabatta}
\index[register]{Bloem}
\index[register]{Poolish}
\index[register]{Gist}

\meta{Ca. 15 uur (incl. nacht) \dvd 1 groot brood \dvd Oven 230\,°C}

\lettrine{C}{iabatta} dankt zijn kenmerkende grote gaten aan twee dingen: een
poolish die een nacht staat te rijzen, en een nat deeg met een hydratatie van
bijna 78\%. Het actieve werk is beperkt, maar reken op geduld: het deeg moet
een nacht en de volgende dag nog zo'n 2,5 uur rijzen.

\blockrule{Ingrediënten}

\noindent\textit{Poolish (voordeeg, avond van tevoren)}
\begin{ingredients}
  \ing{150 g}{tarwebloem}\index[register]{Bloem}
  \ing{150 ml}{water}
  \ing{0,5 g}{gedroogde gist}
\end{ingredients}

\medskip
\noindent\textit{Hoofddeeg}
\begin{ingredients}
  \ing{300 g}{tarwebloem}
  \ing{200 ml}{water}
  \ing{2 g}{gedroogde gist}
  \ing{9 g}{zout (2\%)}
  \ingb{bloem of griesmeel, voor het werkblad}
\end{ingredients}

\blockrule{Bereiding}
\tip{Dit deeg is te nat om te kneden zoals een gewoon brooddeeg. Werk met
natte handen en een royale laag bloem of griesmeel op het werkblad, en gebruik
een deegsteker om het te vormen. Kneed het na de eerste rijs niet meer door,
anders verlies je de luchtbellen waar het deeg zo lang voor heeft staan
rijzen. Neem tarwebloem met een hoog eiwitgehalte, minstens 13\%, zoals
Italiaanse tipo 0 (tarwebloem, geen patentbloem) of Franse T65: gewone bloem houdt het water minder goed
vast en het brood blijft dan plat.}

\medskip
\noindent\textit{Poolish}
\begin{steps}
  \item Meng de bloem, het water en de gist tot een dun beslag. Dek af en laat
        14 uur op kamertemperatuur rijzen, tot het mengsel bubbels heeft en
        naar gist ruikt.
\end{steps}

\medskip
\noindent\textit{Hoofddeeg kneden}
\begin{steps}
  \item Meng de poolish met de bloem, het water en de gist. Kneed, het liefst
        met een keukenmachine, ongeveer 15 minuten tot een soepel maar nog
        plakkerig deeg.
  \item Voeg het zout toe en kneed nog even door tot het is opgenomen.
\end{steps}

\medskip
\noindent\textit{Rijzen en vouwen}
\begin{steps}
  \item Laat het deeg in een ingevette kom 1,5 uur rijzen. Rek en vouw het
        deeg daarbij drie keer, telkens met 30 minuten ertussen: til een rand
        op en vouw hem over het midden, draai de kom een kwartslag en herhaal
        rondom.
  \item Laat het deeg daarna nog 30 minuten met rust in de kom.
\end{steps}

\medskip
\noindent\textit{Vormen en bakken}
\begin{steps}
  \item Stort het deeg voorzichtig op een flink bebloemd werkvlak en vouw het
        losjes tot één langwerpig brood. Leg het op bakpapier, zonder verder
        te kneden of te drukken.
  \item Dek af met een doek en laat 30 minuten rijzen. Verwarm ondertussen de
        oven 30 minuten voor op 230\,°C, het liefst met een bakplaat of
        pizzasteen mee-opwarmend.
  \item Schuif het brood de oven in en voeg stoom toe, bijvoorbeeld met een
        schaaltje water onderin de oven. Bak 25 minuten.
  \item Haal na de eerste 10 minuten het schaaltje water uit de oven, zodat de
        korst knapperig wordt.
  \item Laat de ciabatta afkoelen op een rooster voordat je hem aansnijdt.
\end{steps}

\heroimagefade{images/ciabatta}
\end{recipe}
